% ==============================================================================
% Section 3: Dividend Yield Analysis and the Bhuvanesh Hurdle Criterion
% ==============================================================================

\clearpage
\section[Dividend Yield Analysis]{Dividend Yield Analysis and the Bhuvanesh Hurdle Criterion}

\subsection{Theoretical Premise and the Bhuvanesh Decision Rule}
The \textbf{Dividend Yield} method evaluates the direct annual cash distribution an investor receives relative to the secondary market price paid per share. Under the analytical framework in the course notes (\textit{Bhuvanesh Rule}), dividend yield acts as an immediate income hurdle: if the current dividend yield is less than the investor's opportunity cost of capital ($\Ke = 14.0\%$), the asset fails to compensate for capital risk on immediate cash flow alone and is classified as \textbf{Overvalued}.

\begin{formulabox}[Governing Equation and Bhuvanesh Hurdle Criterion]
\begin{equation}
    \text{Current Yield} = \frac{\text{Proxy DPS}}{P_0} \times 100
\end{equation}
\vspace{-2mm}
\begin{equation}
    \text{Decision Rule}: \quad
    \begin{cases}
        \text{Current Yield} < \Ke \, (14.0\%) \implies \textbf{Overvalued} \\
        \text{Current Yield} \ge \Ke \, (14.0\%) \implies \textbf{Undervalued}
    \end{cases}
\end{equation}
\end{formulabox}

\subsection{Cross-Sectional Dividend Yield Results}
Table~\ref{tab:dividend_yield} details the projected FY26 EPS, standardized $70\%$ proxy DPS, current market prices, computed dividend yields, and comparative verdicts against the $14.0\%$ hurdle rate.

\begin{table}[H]
\centering
\footnotesize
\setlength{\tabcolsep}{2.5pt}
\renewcommand{\arraystretch}{1.08}
\begin{tabular*}{\textwidth}{@{\extracolsep{\fill}} l r r r c c c}
\toprule
\textbf{Company} & \shortstack[r]{\textbf{FY26 EPS}\\\scriptsize(\inr)} & \shortstack[r]{\textbf{Proxy DPS}\\\scriptsize(70\%, \inr)} & \shortstack[r]{\textbf{Market}\\\textbf{Price} (\inr)} & \shortstack[c]{\textbf{Current}\\\textbf{Yield} (\%)} & \shortstack[c]{\textbf{Hurdle ($\bm{\Ke}$)}\\\scriptsize(14.0\%)} & \shortstack[c]{\textbf{Yield}\\\textbf{Verdict}} \\
\midrule
\textbf{Hindustan Unilever}   & 64.08  & 44.86 & 1,973 & 2.27\% & 14.0\% & \textbf{Overvalued} \\
\textbf{ITC Ltd}              & 16.77  & 11.74 & 264   & 4.45\% & 14.0\% & \textbf{Overvalued} \\
\textbf{Nestle India}         & 18.37  & 12.86 & 1,409 & 0.91\% & 14.0\% & \textbf{Overvalued} \\
\textbf{Varun Beverages}      & 4.53   & 3.17  & 204   & 1.55\% & 14.0\% & \textbf{Overvalued} \\
\textbf{Britannia Industries} & 105.71 & 74.00 & 5,120 & 1.44\% & 14.0\% & \textbf{Overvalued} \\
\textbf{Marico Ltd}           & 13.95  & 9.77  & 814   & 1.20\% & 14.0\% & \textbf{Overvalued} \\
\textbf{Tata Consumer}        & 15.63  & 10.94 & 1,010 & 1.08\% & 14.0\% & \textbf{Overvalued} \\
\textbf{Godrej Consumer}      & 14.86  & 10.40 & 881   & 1.18\% & 14.0\% & \textbf{Overvalued} \\
\textbf{Dabur India}          & 10.56  & 7.39  & 382   & 1.93\% & 14.0\% & \textbf{Overvalued} \\
\textbf{Colgate-Palmolive}    & 49.07  & 34.35 & 1,843 & 1.86\% & 14.0\% & \textbf{Overvalued} \\
\bottomrule
\end{tabular*}
\caption{Cross-Sectional Dividend Yield Analysis and Bhuvanesh Hurdle Verification.}
\label{tab:dividend_yield}
\end{table}

\subsection{Analytical Takeaways from the Yield Criterion}
\begin{enumerate}[label=\textbf{\arabic*.}, topsep=2pt, itemsep=1.5pt, parsep=0pt]
    \item \textbf{100\% Failure of Current Income Hurdle:} Under the strict Bhuvanesh criterion, \textbf{all ten companies are classified as Overvalued}. None provide an immediate cash yield remotely approaching the $14.0\%$ benchmark.
    \item \textbf{Yield Dispersion Across the Universe:} \textbf{ITC Ltd} offers the highest dividend yield at \textbf{$4.45\%$}, followed by \textbf{HUL} at \textbf{$2.27\%$}. Conversely, \textbf{Nestle India} generates the lowest yield at \textbf{$0.91\%$}, returning under one rupee per hundred rupees invested.
    \item \textbf{Economic Rationale:} The substantial gap between current yields ($0.9\% - 4.5\%$) and $\Ke$ ($14.0\%$) demonstrates that secondary markets do not treat FMCG shares as quasi-fixed-income instruments; market valuations price in long-term earnings compounding and future dividend growth.
\end{enumerate}
